\chapter{Meritocracy}
\label{ch:meritocracy}

\epigraph{The more a society proclaims its meritocracy, the more its winners can justify their position and the more its losers blame themselves.}{Michael Sandel, \textit{The Tyranny of Merit}}

\noindent
Meritocracy is the belief that outcomes reflect effort and talent.
It is also one of the most heavily subsidized narratives in modern capitalism, functioning as the legitimation mechanism for inequality.

\section{The Meritocratic Belief System}
\label{sec:ch22-system}

\begin{definition}[Meritocratic Regime]
\label{def:meritocratic-regime}
\index{meritocracy}
The \emph{meritocratic regime} is an ideology $\mathcal{P}_M$ containing:
\begin{enumerate}[nosep]
  \item success reflects ability and effort ($p_1$),
  \item opportunity is broadly equal ($p_2$),
  \item failure reflects personal deficiency ($p_3$),
  \item inequality is therefore justified ($p_4$).
\end{enumerate}
These propositions are mutually reinforcing: high credence in $p_1$ raises credence in $p_3$ and $p_4$.
\end{definition}

\section{Institutional Repetition}
\label{sec:ch22-repetition}

The meritocratic regime is reinforced through:
\begin{itemize}[nosep]
  \item educational systems (grading, ranking, credentialing),
  \item corporate hiring and promotion (``talent management''),
  \item media (rags-to-riches narratives, billionaire hagiography),
  \item policy (means-testing, workfare, individual responsibility discourse).
\end{itemize}

Each institution treats meritocratic assumptions as operational truth, creating the conditions under which the assumptions appear self-evident.

\section{Credential Systems as Epistemic Capital Laundering}
\label{sec:ch22-credentials}

\begin{proposition}[Credential Laundering]
\label{prop:credential}
\index{credential laundering}
Credential systems (degrees, certifications, rankings) convert inherited advantage into apparent merit.
If access to credentials correlates with parental wealth ($r > 0$), then the credential $c_j$ for agent $j$ satisfies
\[
  c_j = f(\text{ability}_j, \text{effort}_j) + g(\text{wealth}_j) + \varepsilon_j
\]
The meritocratic narrative claims $g \equiv 0$.
The doxonomic analysis estimates $g$ and reveals the laundering function.
\end{proposition}

\section{A Model of Meritocratic Belief}
\label{sec:ch22-model}

\begin{model}[Meritocratic Belief Production]
\label{mod:meritocracy}
The strength of meritocratic belief $\belief{M}$ is a function of inequality visibility $V$ and narrative investment $\phi_M$:
\[
  \belief{M} = \frac{\phi_M}{\phi_M + \kappa \cdot V}
\]
where $\kappa > 0$ measures the sensitivity of belief to visible inequality.
High investment $\phi_M$ can sustain meritocratic belief even at high inequality, but the required investment grows with $V$.
\end{model}

\section{The Doxonomic Genealogy of Meritocracy}
\label{sec:ch22-genealogy}

The term ``meritocracy'' was coined satirically by Michael Young in \textit{The Rise of the Meritocracy} (1958) as a \emph{warning}, not a recommendation.
Its transformation into a term of praise is itself a doxonomic achievement.

Key institutional channels:
\begin{itemize}[nosep]
  \item \textbf{Standardized testing} (SAT, GRE, IQ tests): created the appearance of objective merit measurement while correlating heavily with parental income and educational investment.
  \item \textbf{Business schools}: propagated the idea that managerial success reflects personal talent, justified executive compensation, and trained a professional class that internalized meritocratic assumptions.
  \item \textbf{Self-help and motivational industry}: annual revenues exceeding \$10 billion in the U.S., reinforcing the belief that success is a matter of mindset and effort.
  \item \textbf{Rags-to-riches media narratives}: disproportionate coverage of exceptional upward mobility, creating availability bias about its frequency.
  \item \textbf{Policy design}: means-testing, workfare, and conditional welfare systems that operationalize the assumption that poverty reflects insufficient effort.
\end{itemize}

\section{Negative Cases: When Meritocratic Belief Fails}
\label{sec:ch22-negative}

Not all investments in meritocratic narrative succeed.
Doxonomic analysis requires attention to failure:
\begin{itemize}[nosep]
  \item In highly visible economic crises (2008, COVID-19), meritocratic belief temporarily declines as structural causes become salient.
  \item In societies with strong labor movements or social-democratic traditions, meritocratic belief is weaker despite similar media exposure.
  \item When inequality becomes extreme enough that the gap between narrative and lived experience is too large, the legitimation function (Definition~\ref{def:legitimation}) breaks down.
\end{itemize}

These negative cases constrain the model: doxonomic investment is necessary but not sufficient for common-sense capture.
Lived experience, counter-institutions, and competing narratives limit the reach of any regime.

\section{Racial and Gender Dimensions}
\label{sec:ch22-racial}

Meritocratic belief is not distributed uniformly.
It intersects with racial and gender regimes (Chapter~\ref{ch:race-coloniality-gender}) in specific ways:
\begin{itemize}[nosep]
  \item White respondents in the U.S.\ consistently score higher on meritocratic belief scales than Black respondents---the people for whom the narrative is most evidently contradicted by structural barriers are least persuaded.
  \item Meritocratic narratives often implicitly assume a subject who is male, able-bodied, and free of caregiving responsibilities.
  \item The ``model minority'' narrative applies meritocratic logic selectively to Asian Americans, using their achievements to delegitimize structural claims by other racialized groups.
\end{itemize}

These patterns confirm that doxonomic capture is partial and contested, not total.

\section{Notes and References}
\label{sec:ch22-notes}

On meritocracy as ideology, see \textcite{sandel2012} and \textcite{bourdieu1984}.
Inequality and its legitimation are analyzed in \textcite{piketty2014} and \textcite{stiglitz2012}.
The role of education in reproducing class structure is central to \textcite{bourdieu1984} and \textcite{bowles1976}.
On credentialism, see Collins (1979).
On the racial dimensions of meritocratic belief, see \textcite{hochschild2016}.
The self-help industry as a doxonomic channel connects to \textcite{ehrenreich2009}.

\section{Exercises}

\begin{exercise}
Using Model~\ref{mod:meritocracy}, compute the narrative investment $\phi_M$ required to maintain $\belief{M} = 0.7$ when inequality doubles.
\end{exercise}

\begin{exercise}
Design a doxometric survey to measure meritocratic belief across socioeconomic groups.
Predict which groups will score highest and explain why doxonomically.
\end{exercise}

\begin{exercise}
Analyze the genre of ``billionaire biography'' as a doxonomic product.
Identify the belief production chain from publisher to reader.
\end{exercise}
